The three triangles joining the interior point to the vertices partition the original triangle. Thus $\boxed{ad_a+bd_b+cd_c=2\Delta}$. If the triangle is equilateral with common side $s$, this gives $d_a+d_b+d_c=\dfrac{2\Delta}{s}$, its altitude. Conversely, if the unweighted sum were constant throughout the interior, continuity would give the same value as the point approaches each vertex. These three limits are the altitudes $\dfrac{2\Delta}{a}$, $\dfrac{2\Delta}{b}$, $\dfrac{2\Delta}{c}$. They can all agree only if $a=b=c$. So the weighted identity is the general theorem, while the unweighted sum is constant exactly for equilateral triangles.
\begin{asy}
import graph;import geometry;usepackage("amsmath");size(7cm);
pair A=(0,3),B=(0,0),C=(5,0),P=(1,1);draw(A--B--C--cycle);draw(P--A^^P--B^^P--C,dashed);draw(P--(0,1)^^P--(1,0),blue);pair Q=(5-5*((5-1)*5+1*3)/34,3*((5-1)*5+1*3)/34);draw(P--Q,blue);label("$A$",A,(0,1));label("$B$",B,(-1,-1));label("$C$",C,(1,-1));label("$P$",P,(0,1));label("$a$",(B+C)/2,(0,-1));label("$b$",(A+C)/2,(1,1));label("$c$",(A+B)/2,(-1,0));
\end{asy}

